\begin{table*}[t]
\centering
\small
\setlength{\tabcolsep}{5pt}
\begin{tabular}{lrrrr}
\toprule
\textbf{Intervention}
& \textbf{WritingBench}
& \textbf{HelloBench}
& \textbf{DoLoMiTes}
& \textbf{Across benchmarks} \\
\midrule
\multicolumn{5}{l}{\textit{Confirmatory single-factor interventions}} \\
\condNoGoals
& \WritingFourFiveRateMean
& \HelloFourFiveRateMean
& \DoLoFourFiveRateMean
& \PooledFourFiveRateMean \\
\condFixedOrder
& \WritingFourSixRateMean
& \HelloFourSixRateMean
& \DoLoFourSixRateMean
& \PooledFourSixRateMean \\
\midrule
\multicolumn{5}{l}{\textit{Diagnostic bundled intervention}} \\
\condSingleWriter
& \WritingFourSevenRateMean
& \HelloFourSevenRateMean
& \DoLoFourSevenRateMean
& \PooledFourSevenRateMean \\
\bottomrule
\end{tabular}
\caption{Pairwise intervention outcomes against \condAgenticCogWriter. Benchmark columns report \condAgenticCogWriter{}'s arithmetic mean of replicate-specific win rates among non-tied comparisons. For the final column, wins, losses, and ties are first pooled across benchmarks within each replication, and the resulting replicate-level win rates are then averaged. These means are descriptive summaries. The first two rows are confirmatory single-factor interventions; \condSingleWriter\ is pre-specified but changes several execution factors together and is interpreted diagnostically. The post-hoc \condSingleContext\ check is reported separately in the text and Appendix~\ref{app:single-context}.}
\label{tab:ablation-results}
\end{table*}
